% Vietnamese translation for OI Public Library
% Source-PDF: IOI中国国家候选队论文/1999/齐鑫.pdf
\documentclass[11pt]{article}
\usepackage{oipl-vn}

\title{Tối ưu hóa cắt tỉa trong phương pháp tìm kiếm}
\author{Qi Xin}
\date{}

\begin{document}
\maketitle

\origtitle{Pruning optimization in search methods}
\sourcepath{IOI China candidate team papers / 1999 / Qi Xin.pdf}

\paragraph{Từ khóa.}
Tìm kiếm; tối ưu hóa; cắt tỉa.

\section*{Tóm tắt}

Bài viết này thảo luận một kỹ thuật tối ưu hóa thường gặp nhất trong phương
pháp tìm kiếm: cắt tỉa. Trọng tâm chính là thiết kế tiêu chuẩn phán đoán để
cắt tỉa. Trước hết, bài viết dùng cây tìm kiếm để giải thích một cách trực
quan cắt tỉa là gì; sau đó phân tích ba nguyên tắc khi thiết kế tiêu chuẩn
cắt tỉa: đúng đắn, chính xác và hiệu quả. Bài viết chia các hướng thiết kế
tiêu chuẩn cắt tỉa thường gặp thành hai loại lớn, cắt tỉa theo tính khả thi
và cắt tỉa theo tính tối ưu, rồi dùng mỗi loại một bài thi làm ví dụ để minh
họa ba nguyên tắc trên. Cuối cùng, bài viết tổng kết một số nhận xét về
phương pháp cắt tỉa.

\section{Dẫn nhập}

Tìm kiếm là một phương pháp cơ bản trong trí tuệ nhân tạo, đồng thời cũng là
một phương pháp mà thí sinh thi tin học phải nắm vững. Khi xây dựng một thuật
toán tìm kiếm, hai vấn đề đầu tiên thường không nằm ngoài:
\begin{enumerate}
  \item Xây dựng cấu trúc thuật toán.
  \item Chọn cấu trúc dữ liệu thích hợp.
\end{enumerate}

Tuy nhiên, như mọi người đều biết, độ phức tạp thời gian của phần lớn phương
pháp tìm kiếm là cấp mũ. Một tìm kiếm đơn giản không tối ưu hóa thường có
hiệu suất thời gian thấp đến mức khó chấp nhận, càng khó đáp ứng giới hạn
thời gian nghiêm ngặt trong thi tin học.

Nội dung chính của bài viết này là một phương pháp cơ bản để tối ưu hóa
chương trình sau khi cấu trúc thuật toán đã được xây dựng: cắt tỉa.

Trước hết cần làm rõ ý nghĩa của ``cắt tỉa''. Ta biết rằng quá trình tìm kiếm
có thể xem như quá trình xuất phát từ gốc cây và duyệt một cây lộn ngược,
tức cây tìm kiếm. Gọi là cắt tỉa vì thông qua một tiêu chuẩn phán đoán nào
đó, ta tránh được một số quá trình duyệt không cần thiết. Nói hình ảnh hơn,
ta cắt bỏ một số ``cành'' trong cây tìm kiếm, nên gọi là cắt tỉa.

\begin{center}
\begin{tabular}{c}
\(\bullet\)\\[-0.2em]
\(/\quad|\quad\backslash\)\\[-0.2em]
\(\bullet\quad\bullet\quad\bullet\)\\[-0.2em]
\(/\ \backslash\qquad /\ \backslash\qquad /\ \backslash\)\\[-0.2em]
\(\bullet\ \bullet\quad\bullet\ \bullet\quad\bullet\ \bullet\)
\end{tabular}

\smallskip
Hình 1. Cây tìm kiếm, trong đó một số nhánh có thể bị loại bỏ trước khi duyệt.
\end{center}

Khi viết chương trình tìm kiếm, nói chung ta đều phải xét đến cắt tỉa. Hiển
nhiên, vấn đề cốt lõi của tối ưu hóa bằng cắt tỉa là thiết kế tiêu chuẩn
phán đoán: phương pháp xác định nhánh nào nên bỏ, nhánh nào nên giữ. Một tiêu
chuẩn tốt thường có thể rút ngắn đáng kể thời gian chạy; ngược lại, nó cũng
có thể gây phản tác dụng. Vì vậy trước hết ta cần phân tích các nguyên tắc
cần tuân theo khi thiết kế tiêu chuẩn cắt tỉa.

\section{Các nguyên tắc cắt tỉa}

\paragraph{Nguyên tắc thứ nhất: đúng đắn.}
Như đã nói ở trên, cắt tỉa có thể tối ưu hóa hiệu suất thực thi vì nó có thể
``cắt bỏ'' một số ``cành'' của cây tìm kiếm. Nhưng nếu khi cắt tỉa ta cắt luôn
nhánh ``mọc ra'' lời giải cần tìm, mọi tối ưu hóa đều mất ý nghĩa. Vì vậy yêu
cầu đầu tiên đối với cắt tỉa là tính đúng đắn: phải bảo đảm không đánh mất
kết quả đúng. Đây là tiền đề của tối ưu hóa bằng cắt tỉa.

Để thỏa mãn nguyên tắc này, ta nên dùng \term{điều kiện cần} để phán đoán cắt
tỉa. Nói cách khác, ta xem nhánh đang xét có thể bị cắt hay không thông qua
các đặc trưng mà lời giải bắt buộc phải có, các điều kiện mà lời giải bắt
buộc phải thỏa, v.v. Như vậy có thể bảo đảm rằng nhánh bị cắt chắc chắn
không phải nhánh chứa lời giải đúng. Tất nhiên, theo định nghĩa điều kiện cần,
không bị cắt không có nghĩa là nhất định sẽ thu được lời giải đúng, nếu không
thì cũng chẳng cần tìm kiếm nữa.

\paragraph{Nguyên tắc thứ hai: chính xác.}
Trên cơ sở bảo đảm đúng đắn, yêu cầu thứ hai đối với tiêu chuẩn cắt tỉa là
tính chính xác, tức là có thể cắt càng nhiều nhánh không dẫn đến lời giải
đúng càng tốt. Chỉ khi có độ chính xác cao, cắt tỉa mới thật sự phát huy hiệu
quả tối ưu hóa. Vì thế có thể nói chính xác là sinh mệnh của tối ưu hóa bằng
cắt tỉa.

Dĩ nhiên, để nâng cao độ chính xác của tiêu chuẩn cắt tỉa, ta phải phân tích
đặc điểm của bài toán một cách toàn diện và tỉ mỉ, cố gắng phát hiện bản chất
của bài toán để thiết kế được phương án phán đoán tốt.

\paragraph{Nguyên tắc thứ ba: hiệu quả.}
Nói chung, sau khi thiết kế tiêu chuẩn cắt tỉa, ta phải thực hiện một phép
phán đoán trên mỗi nhánh của cây tìm kiếm. Nhưng vì phép phán đoán dựa trên
``điều kiện cần'' để có lời giải, chắc chắn vẫn có nhiều nhánh không chứa lời
giải đúng mà không bị cắt. Trong các trường hợp ấy, thao tác phán đoán cắt
tỉa rõ ràng có tác dụng phụ đối với hiệu suất chương trình. Để giảm tác dụng
phụ này, ngoài việc cải thiện độ chính xác của phán đoán, ta thường còn phải
nâng cao hiệu suất thời gian của chính thao tác phán đoán.

Điều này dẫn đến một mâu thuẫn: để tăng hiệu quả tối ưu hóa, ta phải nâng cao
độ chính xác của tiêu chuẩn cắt tỉa, nên thường buộc phải tăng độ phức tạp
của thao tác phán đoán, đồng thời làm giảm hiệu suất thời gian của nó. Nhưng
nếu thời gian tiêu tốn cho phán đoán quá lớn, nó có thể làm giảm, thậm chí
triệt tiêu hoàn toàn hiệu quả tối ưu hóa do độ chính xác cao hơn mang lại.
Trong nhiều trường hợp, việc giải quyết tốt mâu thuẫn này thường trở thành
điểm then chốt của tối ưu hóa thuật toán tìm kiếm.

Tóm lại, các nguyên tắc chính của tối ưu hóa bằng cắt tỉa có thể quy về ba
từ: đúng, chính xác, hiệu quả.

Khi áp dụng cắt tỉa, chỉ có các nguyên tắc trên vẫn chưa đủ; ta còn cần
nghiên cứu cụ thể các hướng thiết kế tiêu chuẩn cắt tỉa. Có thể chia đại khái
các phán đoán cắt tỉa thường dùng thành hai loại:
\begin{enumerate}
  \item Cắt tỉa theo tính khả thi.
  \item Cắt tỉa theo tính tối ưu, hay cắt tỉa bằng cận trên và cận dưới.
\end{enumerate}
Sau đây ta sẽ kết hợp ba nguyên tắc trên để thảo luận từng loại.

\section{Cắt tỉa theo tính khả thi}

Ta đã biết quá trình tìm kiếm có thể xem như việc duyệt một cây. Trong rất
nhiều trường hợp, không phải mọi nhánh của cây tìm kiếm đều có thể dẫn đến
kết quả ta cần; nhiều nhánh thực chất chỉ là ngõ cụt. Nếu ngay khi vừa bước
vào một ngõ cụt như vậy ta đã nhận ra và lập tức cắt tỉa, hiệu suất chương
trình thường sẽ được cải thiện. Cắt tỉa theo tính khả thi chính là xuất phát
từ suy nghĩ này.

Hãy xét một ví dụ: ``Betsy đi du lịch'' của USACO.

\paragraph{Tóm tắt đề bài.}
Một thị trấn hình vuông được chia thành \(N^2\) ô vuông nhỏ. Betsy cần đi từ
ô góc trái trên đến ô góc trái dưới, và đi qua mỗi ô đúng một lần. Với \(N\)
cho trước, hãy tính số tuyến đường mà Betsy có thể đi.

Ta dùng phương pháp quay lui theo chiều sâu để giải bài này. Betsy xuất phát
từ góc trái trên; mỗi bước có thể di chuyển từ một ô sang một ô kề cạnh chưa
đi qua. Gặp ngõ cụt thì quay lui. Khi đã đi \(N^2-1\) bước và đến góc trái
dưới, ta thu được một đường mới, rồi tiếp tục quay lui tìm kiếm cho đến khi
duyệt hết mọi đường.

Tuy nhiên, nếu chỉ lập trình theo khung thuật toán trên, hiệu suất thời gian
rất thấp; trường hợp \(N=6\) cũng không được xử lý tốt. Vì vậy tối ưu hóa là
điều bắt buộc. Điểm then chốt của tối ưu hóa bài này là ở một bước tìm kiếm
nào đó, có thể phán đoán trước liệu trong phần tìm kiếm phía sau chắc chắn sẽ
gặp ngõ cụt hay không. Cắt tỉa theo tính khả thi có thể phát huy tác dụng tại
đây.

Trước hết ta phân tích phương pháp cắt tỉa từ góc độ ``điều kiện cần'', tức
các đặc trưng mà một lời giải hợp lệ phải có. Có hai hướng chính:
\begin{enumerate}
  \item Với một đường đi hợp lệ, ngoài điểm xuất phát và ô đích, mỗi ô trung
  gian tất yếu có quá trình ``vào một lần, ra một lần''. Vì vậy trong quá
  trình tìm kiếm phải bảo đảm mỗi ô chưa đi qua còn kề với ít nhất hai ô chưa
  đi qua, trừ khi lúc ấy Betsy đang đứng ngay bên cạnh nó. Đây là phân tích ở
  góc độ vi mô.

  \item Trên cơ sở điều kiện thứ nhất, ta còn có thể phân tích từ góc độ vĩ
  mô để nâng cao thêm độ chính xác của cắt tỉa. Rõ ràng tại mọi thời điểm của
  một phương án di chuyển hợp lệ, không thể tồn tại vùng bị cô lập. Dù mỗi ô
  trong vùng cô lập vẫn có thể có ít nhất hai ô trống kề cạnh, xét cả vùng
  như một chỉnh thể thì Betsy đã không thể đi tới đó nữa. Ta cũng nên kịp
  thời nhận ra và tránh trường hợp này.
\end{enumerate}

Hai điều kiện cắt tỉa trên đều đúng và có độ chính xác khá cao. Nhưng chỉ
thỏa hai điểm ấy vẫn chưa đủ; quá trình phán đoán cắt tỉa còn phải cố gắng
hiệu quả. Nếu mỗi lần phán đoán ta đơn giản quét toàn bộ \(N^2\) ô, hiệu suất
thấp là điều dễ thấy. Do đó phải đơn giản hóa quá trình phán đoán hết mức có
thể.

Thực ra, mỗi lần Betsy di chuyển chỉ ảnh hưởng đến các ô gần đó. Vì vậy mỗi
lần thực hiện phán đoán cắt tỉa, ta chỉ nên kiểm tra các ô lân cận cô ấy.

Với điều kiện cắt tỉa thứ nhất, có thể đặt một mảng đánh dấu kiểu nguyên, lưu
riêng cho mỗi ô số ô kề cạnh chưa được đi qua. Mỗi lần Betsy đi đến vị trí
mới, nhiều nhất chỉ có bốn ô kề cạnh vị trí ấy bị thay đổi giá trị đánh dấu,
nên chỉ cần kiểm tra các giá trị này.

Với điều kiện cắt tỉa thứ hai, xử lý hơi phức tạp hơn. Nhưng ta vẫn có thể
dùng phương pháp phân tích cục bộ: chỉ dựa vào các ô gần Betsy để xác định có
nên cắt tỉa hay không. Nguyên lý có thể mô tả như sau. Ký hiệu \(\mathrm{X}\)
là ô đã đi qua, \(\mathrm{O}\) là ô chưa đi qua, \(B\) là vị trí mới của
Betsy, và mũi tên là hướng di chuyển:
\[
\begin{array}{ccc@{\qquad}ccc@{\qquad}ccc}
  \mathrm{O} & \mathrm{X} & \mathrm{O} &
  \mathrm{X} & \mathrm{X} & \mathrm{O} &
  \mathrm{O} & \mathrm{X} & \mathrm{X} \\
  \mathrm{O} & B\!\uparrow & \mathrm{O} &
  \mathrm{O} & B\!\uparrow & \mathrm{O} &
  \mathrm{O} & B\!\uparrow & \mathrm{O}
\end{array}
\]
Ba cấu hình trên là ba tình huống có thể cắt tỉa. Vì mọi ô Betsy đã đi qua
chắc chắn tạo thành một vùng liên thông bốn hướng, trong mỗi tình huống hai ô
trắng được tách bởi vệt đã đi qua chắc chắn không liên thông với nhau. Ít nhất
một trong hai phía thuộc về một vùng cô lập, nên đều có thể cắt tỉa.

Sau các tối ưu hóa trên, hiệu suất thời gian của chương trình được cải thiện
rất nhiều; xem phụ lục.

Nói chung, cắt tỉa theo tính khả thi thường dùng trong các bài tìm kiếm đường
đi. Ngoài ví dụ này, các bài như ``Prime Circle'' ở ACM Asian Regional 1996
cũng có thể dùng phương pháp cắt tỉa này.

Khi áp dụng cắt tỉa theo tính khả thi, trước hết cần phân tích đặc điểm bài
toán từ nhiều góc độ, tìm càng nhiều tình huống có thể cắt tỉa càng tốt. Đồng
thời, cũng phải chú ý nâng cao hiệu suất thời gian của cắt tỉa, nên ở đây ta
dùng phương pháp ``phán đoán cục bộ''. Đặc biệt khi xử lý điều kiện cắt tỉa
thứ hai, ta dùng phán đoán cục bộ để thể hiện tính chất toàn cục, tức có tồn
tại vùng cô lập hay không. Kỹ thuật này không chỉ có ích khi thiết kế phương
pháp cắt tỉa, mà còn được ứng dụng rất rộng rãi ở nhiều mặt khác.

\section{Cắt tỉa theo tính tối ưu}

Trong các bài ta thường gặp, có một lớp lớn gọi là bài toán tối ưu hóa, tức
kết quả cần tìm là lời giải tối ưu. Nếu dùng tìm kiếm để giải các bài như vậy,
cắt tỉa theo tính tối ưu là điều chắc chắn phải xét đến.

Để thống nhất cách trình bày, trước hết cần nói rõ một điểm. Ta biết rằng độ
tốt xấu của lời giải thường được đánh giá bằng một hàm đánh giá. Ở đây định
nghĩa một hàm đánh giá trừu tượng gọi là ``độ ưu'', giá trị càng lớn thì lời
giải tương ứng càng tốt. Với bài cụ thể, có thể xem ``độ ưu'' là lợi ích
dương, hoặc là chi phí lấy dấu âm, v.v.

Sau đó, hãy nhìn lại quá trình tìm kiếm lời giải tối ưu. Nói chung, ta cần
lưu một ``lời giải tốt nhất hiện tại'', thực chất là lưu một cận dưới của độ
ưu. Khi duyệt đến nút lá của cây tìm kiếm, ta thu được một lời giải mới và do
đó cũng thu được giá trị hàm đánh giá của nó. So sánh giá trị ấy với cận dưới
đang lưu; nếu độ ưu của lời giải mới lớn hơn, nó trở thành cận dưới mới. Sau
khi tìm kiếm kết thúc, lời giải đang lưu là lời giải tối ưu.

Vậy cắt tỉa theo tính tối ưu tiến hành như thế nào? Khi đang ở trên một nhánh
của cây tìm kiếm, ta có thể dùng một phương pháp nào đó để ước lượng cận trên
của hàm đánh giá đối với mọi lời giải trên nhánh ấy, tức hàm ước lượng \(h\).
Hiển nhiên, \(h\) lớn hơn cận dưới độ ưu hiện tại là điều kiện cần để trên
nhánh ấy tồn tại lời giải tối ưu; nếu không, chắc chắn có thể cắt tỉa. Vì vậy
cắt tỉa theo tính tối ưu cũng có thể gọi là cắt tỉa bằng cận trên và cận
dưới. Đồng thời ta cũng thấy vấn đề cốt lõi của loại cắt tỉa này là xây dựng
hàm ước lượng.

Sau đây là một ví dụ điển hình: ``số lần nhân ít nhất''.

\paragraph{Tóm tắt đề bài.}
Bắt đầu từ \(x\), hãy dùng số phép nhân ít nhất để thu được \(x^n\), trong đó
\(n\) là dữ liệu vào.

Vì nhân hai lũy thừa tương đương cộng hai số mũ, bài này có thể biểu diễn
tương đương như sau. Cần xây dựng một dãy \(\{a_i\}\) thỏa
\[
  a_i =
  \begin{cases}
    1, & i=1,\\
    a_j+a_k, & 1\le j,k<i,\ i>1,
  \end{cases}
\]
yêu cầu \(a_t=n\) và làm cho \(t\) nhỏ nhất.

Ta chọn quay lui làm cấu trúc chính. Khi tìm kiếm đến tầng \(i\), giá trị
\(a_i\) nằm trong khoảng từ \(a_{i-1}+1\) đến \(2a_{i-1}\). Để nhanh chóng
tiến gần mục tiêu \(n\), nên thử giá trị cho \(a_i\) từ \(2a_{i-1}\) giảm
dần, dĩ nhiên các giá trị được chọn không được vượt quá \(n\). Khi tìm thấy
\(n\), ta thu được một lời giải và so sánh với lời giải đang lưu để lấy tốt
hơn. Khi tìm kiếm kết thúc, ta thu được lời giải tối ưu cuối cùng.

Nếu trực tiếp tìm kiếm theo cấu trúc trên, hiệu suất rõ ràng rất thấp. Vì vậy
ta có thể dùng cắt tỉa theo tính tối ưu để tối ưu hóa.

\paragraph{Tối ưu thứ nhất.}
Trước hết phải xây dựng hàm ước lượng. Vì nhân đôi số lớn nhất trong dãy chắc
chắn là cách tăng nhanh nhất, một hàm ước lượng đơn giản nhất là, giả sử số
lớn nhất hiện tại trong dãy là \(a_i\), tức độ sâu tìm kiếm hiện tại là \(i\):
\[
  h=\left\lceil \log_2 \frac{n}{a_i}\right\rceil .
\]

Tuy nhiên, độ chính xác của hàm ước lượng này quá kém. Đặc biệt khi \(a_i\)
lớn hơn \(n/2\), \(h\) chỉ bằng \(1\), gần như không thể phát huy tác dụng cắt
tỉa. Vì vậy, dù dùng quay lui theo chiều sâu hay tìm kiếm \(A^\ast\) theo
chiều rộng, nếu vẫn dùng hàm ước lượng này thì hiệu quả tối ưu hóa đều khá
hạn chế.

\paragraph{Tối ưu thứ hai.}
Hướng vào ``sinh mệnh'' của hàm ước lượng, tức độ chính xác, ta có thể dùng
một mở rộng của tính chất cộng theo chẵn lẻ để cải thiện độ chính xác của hàm
ước lượng trong một số trường hợp. Chi tiết cải tiến được nêu trong phụ lục.

\paragraph{Tối ưu thứ ba.}
Hàm ước lượng mới có độ chính xác cao hơn một chút, nhưng độ phức tạp thời
gian cũng tăng tương ứng. Để nâng cao hiệu suất cắt tỉa, ta có thể dùng kỹ
thuật ``tinh chỉnh từng bước'': trước hết dùng hàm ước lượng cũ, nhanh nhưng
kém chính xác, để lọc một lần; sau đó dùng hàm ước lượng mới, chính xác hơn
một chút nhưng chậm hơn, để giảm số trường hợp lọt qua, nhằm cân bằng giữa độ
chính xác và hiệu suất chạy.

\paragraph{Tối ưu thứ tư.}
Trước khi tìm kiếm, ta có thể phân tích \(n\) thành thừa số nguyên tố, rồi
dùng phương pháp tìm kiếm nói trên giải riêng cho từng thừa số nguyên tố. Nếu
một thừa số nguyên tố vẫn quá lớn, còn có thể xét nó trừ đi \(1\) rồi phân
tích tiếp. Điều này tương đương chia quá trình xây dựng dãy thành nhiều giai
đoạn để xử lý. Kết quả thu được tuy chưa chắc là tối ưu toàn cục, nhưng có
thể dùng làm cận dưới ban đầu của độ ưu, giúp quá trình tìm kiếm toàn cục
phía sau tránh rất nhiều đường vòng. Thực tế, phương pháp ước lượng này khá
chính xác; trong nhiều trường hợp nó trực tiếp cho lời giải tối ưu, nhờ đó
hiệu suất chương trình tăng rất mạnh.

\paragraph{Tối ưu thứ năm.}
Khi số lớn nhất hiện tại \(a_i\) trong dãy vượt quá \(n/2\), hàm ước lượng cũ
khó phát huy tác dụng. Nhưng phương án tối ưu sau đó thực chất tương đương
với việc chọn càng ít số càng tốt từ \(i\) số \(a_1,\ldots,a_i\), cho phép lặp
lại, sao cho tổng của chúng bằng \(n\). Bài toán này đã có thể giải bằng quy
hoạch động. ``Hàm ước lượng'' thu được theo cách này không chỉ hoàn toàn
chính xác, mà thậm chí còn có thể thay thế trực tiếp quá trình tìm kiếm phía
sau. Ở đây cũng thể hiện sự bổ sung ưu thế giữa tìm kiếm và quy hoạch động.

Bài này chứa nhiều kỹ thuật cắt tỉa theo tính tối ưu, và hiệu quả tối ưu hóa
của chúng cũng rất rõ rệt; xem phân tích trong phụ lục.

Từ bài này, kết hợp với kinh nghiệm trước đây, ta có thể rút ra một số điểm
cần chú ý khi áp dụng cắt tỉa theo tính tối ưu:
\begin{enumerate}
  \item Thiết kế hàm ước lượng trước hết tất nhiên phải thỏa nguyên tắc đúng
  đắn, tức dùng ``điều kiện cần'' để cắt tỉa. Trên cơ sở đó cần chú ý nâng
  cao độ chính xác của ước lượng. Tối ưu thứ hai của bài này nhằm mục đích
  ấy.

  \item Cũng như các thao tác phán đoán cắt tỉa khác, khi tăng độ chính xác
  của hàm ước lượng trong cắt tỉa theo tính tối ưu, đồng thời phải chú ý làm
  cho quá trình tính toán hiệu quả nhất có thể. Vì phán đoán cắt tỉa được
  thực hiện cực kỳ thường xuyên khi chạy, việc mài giũa thuật toán của nó là
  rất cần thiết. Đôi khi còn có thể dùng một số thủ pháp đặc biệt, như kỹ
  thuật ``tinh chỉnh từng bước'' ở trên, để tìm điểm cân bằng giữa hiệu suất
  thời gian và độ chính xác.

  \item Khi dùng cắt tỉa theo tính tối ưu, một cận dưới ``độ ưu'' ban đầu tốt
  thường rất quan trọng. Trước khi bắt đầu tìm kiếm, ta có thể dùng một
  phương pháp hiệu quả nào đó, chẳng hạn tham lam, để tìm một lời giải tương
  đối tốt làm cận dưới ban đầu; cách này thường cắt bỏ được rất nhiều nhánh
  rõ ràng không thể tốt. Bài này dùng tìm kiếm theo giai đoạn để định cận ban
  đầu, và nhờ đó đạt tối ưu hóa lớn.

  \item Ví dụ trong mục này là việc áp dụng cắt tỉa theo tính tối ưu trong
  tìm kiếm theo chiều sâu. Dĩ nhiên, trong tìm kiếm theo chiều rộng cũng có
  thể dùng cắt tỉa theo tính tối ưu; đó chính là phương pháp nhánh cận thường
  nói. Bài này cũng có thể giải bằng cách kết hợp nhánh cận với thuật toán
  \(A^\ast\), dùng hàm ước lượng cải tiến làm ước lượng cận trên độ ưu, tức
  hàm \(h^\ast\), và dùng quy hoạch động nói trên làm ước lượng cận dưới độ
  ưu. Tuy vậy, hiệu suất thời gian và không gian đều kém hơn. Nhưng trong một
  số bài, sự kết hợp giữa nhánh cận và \(A^\ast\), nhờ tính ổn định khá tốt,
  vẫn có đất dụng võ.
\end{enumerate}

\section{Tổng kết}

Do phương pháp tìm kiếm có ``khiếm khuyết bẩm sinh'' về hiệu suất thời gian,
người ta đã nghiên cứu nhiều kỹ thuật tối ưu hóa nhắm vào nó. ``Cắt tỉa'' mà
bài viết này bàn đến là một trong các phương pháp tối ưu hóa thường gặp nhất;
có thể nói hầu như chỉ cần muốn dùng thuật toán tìm kiếm để giải bài, ta đều
phải xét đến tối ưu hóa bằng cắt tỉa.

Ngoài ra cần nói rõ rằng hai loại phán đoán cắt tỉa được giới thiệu trong bài
này, tính khả thi và tính tối ưu, chỉ được phân loại theo đối tượng ứng dụng
khác nhau. Ba nguyên tắc đã trình bày ở trên, đúng đắn, chính xác và hiệu
quả, mới là linh hồn xuyên suốt từ đầu đến cuối.

Bài viết còn giới thiệu một số kỹ thuật thường dùng trong cắt tỉa, như
``phân tích cục bộ'', ``tinh chỉnh từng bước'', v.v. Dùng chúng thích hợp cũng
có thể nâng cao đáng kể hiệu suất chương trình, chủ yếu là hiệu suất thời
gian.

Nhưng dù phương pháp cắt tỉa tinh xảo đến đâu, nó cũng không thể hạ độ phức
tạp thời gian của thuật toán tìm kiếm về bản chất; đây là sự thật không thể
phủ nhận. Vì vậy, trước khi bắt tay thiết kế một thuật toán tìm kiếm, ta cũng
nên cân nhắc xem có tồn tại phương pháp hiệu quả hơn hay không.

Hơn nữa, trong thi tin học còn có vấn đề độ phức tạp lập trình. Nếu ta dùng
rất nhiều kỹ thuật tối ưu hóa cho một thuật toán tìm kiếm, tuy có thể cải
thiện phần nào hiệu suất thời gian, nhưng thường phải tiêu tốn rất nhiều thời
gian lập trình, rất dễ rơi vào kết quả ``nhặt hạt vừng, đánh rơi quả dưa''.

Tóm lại, trong quá trình thiết kế chương trình thực tế, ta không nên chui vào
ngõ hẹp, mà càng phải phát huy đầy đủ tính linh hoạt của tư duy, kiên trì
phương pháp tư tưởng ``phân tích cụ thể từng vấn đề cụ thể''.

\section*{Tài liệu tham khảo}

\begin{enumerate}
  \item Liu Fusheng, Wang Jiande, \textit{Hướng dẫn Olympic Tin học quốc tế
  thanh thiếu niên: tìm kiếm trí tuệ nhân tạo và thiết kế chương trình}.
  \item \textit{Informatics Olympic}.
\end{enumerate}

\appendix
\section{Phụ lục}

\subsection{Tình hình kiểm thử của ``Betsy đi du lịch''}

Đơn vị thời gian là giây; môi trường kiểm thử là Pentium II 266 MHz, 64 MB.
\[
\begin{array}{c|r|c|c|c|c}
N & \text{Số đường} & \text{Không cắt} & \text{Cắt 1} &
\text{Cắt 2} & \text{Cả hai}\\
\hline
2 & 1 & <0.01 & <0.01 & <0.01 & <0.01\\
3 & 2 & <0.01 & <0.01 & <0.01 & <0.01\\
4 & 8 & <0.01 & <0.01 & <0.01 & <0.01\\
5 & 86 & 0.17 & <0.01 & <0.01 & <0.01\\
6 & 1770 & 41.25 & 0.44 & 0.33 & 0.11\\
7 & 88418 & \text{Rất lâu} & 28.06 & 26.86 & 8.51
\end{array}
\]

Có thể thấy việc dùng các tiêu chuẩn cắt tỉa theo tính khả thi nói trên làm
hiệu suất thời gian của chương trình tăng rất đáng kể.

\subsection{Cải tiến hàm ước lượng của ``số lần nhân ít nhất''}

Ý tưởng thiết kế hàm ước lượng ban đầu thực chất là trên cơ sở dãy hiện tại
\(a_1,\ldots,a_i\), lý tưởng hóa việc xây dựng
\[
  2a_i,4a_i,\ldots,2^p a_i.
\]
Khi \(2^p a_i<n<2^{p+1}a_i\), phương pháp ước lượng cũ cho rằng chỉ cần thêm
một phép cộng nữa, tức tìm một số cộng với \(2^p a_i\), là đủ.

Nhưng điều này chỉ bảo đảm có thể thu được một số lớn hơn hoặc bằng \(n\),
chưa chắc thu được đúng \(n\). Ta phân tích như sau.

Trước hết, mọi số tự nhiên \(i\) đều có thể viết dưới dạng
\[
  2^k(2m+1)\qquad (k,m\in\mathbb{Z}_{\ge 0}).
\]
Đặt \(k(i)\) là giá trị \(k\) tương ứng với số tự nhiên \(i\). Rõ ràng với mọi
hai số tự nhiên \(a,b\), ta có
\[
  k(a+b)\ge \min\{k(a),k(b)\}.
\]
Từ đây dễ liên tưởng đến các tính chất chẵn lẻ như lẻ cộng lẻ bằng chẵn, chẵn
cộng chẵn bằng chẵn, lẻ cộng chẵn bằng lẻ.

Xét dãy được xây dựng trong ước lượng ở trên:
\[
  a_1,a_2,\ldots,a_i,2a_i,4a_i,\ldots,2^p a_i,
  \qquad 2^p a_i<n<2^{p+1}a_i.
\]
Trước khi dùng phán đoán cắt tỉa mới, ta nên kiểm tra
\[
  \left\lceil \log_2\frac{n}{a_i+a_{i-1}}\right\rceil \ge p,
\]
vì điều kiện này bảo đảm rằng chỉ có xây dựng dãy trên mới là tối ưu.

Nếu tồn tại số tự nhiên \(j\), \(1\le j\le p\), sao cho
\[
  k(2^j a_i)>k(n),
\]
thì từ \(k(a+b)\ge\min\{k(a),k(b)\}\), với \(j\le t\le p\) ta có
\[
  k(2^t a_i+2^p a_i)\ge k(2^t a_i)\ge k(2^j a_i)>k(n).
\]
Do \(k(a)=k(b)\) là điều kiện cần để \(a=b\), suy ra
\[
  2^t a_i+2^p a_i\ne n.
\]
Nghĩa là mọi số trong \(2^j a_i,\ldots,2^p a_i\) cộng với \(2^p a_i\) đều
không thể thu được \(n\). Vì vậy, nếu
\[
  n-2^p a_i > 2^{j-1}a_i,
\]
thì sau khi thu được \(2^p a_i\), ít nhất còn cần thêm hai phép cộng nữa mới
có thể thu được \(n\).

Cải tiến trên có thể làm hàm ước lượng chính xác hơn một chút trong một số
trường hợp, nhưng bản thân nó có hiệu suất thời gian kém hơn. Lý do là hàm
ước lượng mới không chỉ bao gồm hàm ước lượng cũ, tức quá trình xây dựng dãy,
mà còn thêm các thao tác như tính \(k(a_i)\). Vì vậy, dù hàm mới chỉ là một
cải tiến của hàm cũ, không phải hai mặt bổ sung cho nhau như trong bài
``Betsy đi du lịch'', trong chương trình thực tế vẫn cần gọi liên tiếp hai
hàm ước lượng, tức dùng kỹ thuật ``tinh chỉnh từng bước'' đã nói ở trên.

Cuối cùng cần nói thêm rằng mọi phần dùng phép toán logarit trong các hàm ước
lượng nói tới trong bài này đều phải đổi thành vòng lặp cộng dồn khi lập
trình thực tế. Nguyên nhân là hàm số thực trong máy tính được tính bằng khai
triển chuỗi, thực chất cũng là cấu trúc lặp, lại phức tạp và kém hiệu quả
hơn. Dĩ nhiên, trong chương trình thực tế còn có một số tối ưu hóa chi tiết
khác.

\subsection{Kiểm thử và phân tích ``số lần nhân ít nhất''}

Một số dữ liệu kiểm thử, đơn vị giây. Môi trường kiểm thử là Pentium Pro
233 MHz, 64 MB.
\[
\begin{array}{r|r|c|c|c|c|c|c|c}
N & \text{Số nhân} & P1 & P2 & P3 & P4 & P5 & P6 & P7\\
\hline
127 & 10 & 0.22 & 0.22 & 0.05 & 0.06 & 0.06 & <0.01 & 0.16\\
191 & 11 & 3.52 & 3.90 & 0.83 & 0.55 & 0.55 & 0.27 & 1.75\\
382 & 11 & 13.02 & 13.29 & 1.05 & 0.94 & 0.66 & 0.33 & 2.69\\
511 & 12 & \text{Rất lâu} & \text{Rất lâu} & 0.82 & 5.66 & 0.33 & 0.27 & 1.81\\
635 & 13 & 46.96 & 40.37 & 20.65 & 19.61 & 10.06 & 8.13 & \text{Tràn}\\
719 & 13 & 45.59 & 39.11 & 45.64 & 13.45 & 39.11 & 6.64 & 27.35\\
1002 & 13 & 11.81 & 10.65 & 2.74 & 3.46 & 1.10 & 0.99 & 6.26\\
1023 & 13 & \text{Rất lâu} & \text{Rất lâu} & 2.66 & \text{Rất lâu} & 1.04 & 0.88 & 6.15\\
1357 & 13 & 9.51 & 6.98 & 5.93 & 5.44 & 3.40 & 2.74 & 13.35\\
1894 & 14 & 44.87 & 37.24 & 13.89 & 15.49 & 6.37 & 4.61 & 23.68
\end{array}
\]

Khi \(n\le 1000\), thời gian chạy của chương trình 6 đều không vượt quá
10 giây, trong đó chỉ có hơn mười dữ liệu kiểm thử chạy trên 5 giây.

\[
\begin{array}{c|c|c|c|c|c}
\text{Chương trình} & \text{Cắt đơn giản} & \text{Cắt cải tiến} &
\text{Định cận đầu} & \text{Quy hoạch động} & A^\ast+\text{nhánh cận}\\
\hline
1 & \text{có} & & & &\\
2 & \text{có} & \text{có} & & &\\
3 & \text{có} & & \text{có} & &\\
4 & \text{có} & & & \text{có} &\\
5 & \text{có} & \text{có} & \text{có} & &\\
6 & \text{có} & \text{có} & \text{có} & \text{có} &\\
7 & & & & & \text{có}
\end{array}
\]

Trước hết so sánh hiệu quả tối ưu hóa khi ba điểm tối ưu chính được dùng
riêng lẻ: phán đoán cắt tỉa cải tiến, tức tinh chỉnh từng bước; định cận ban
đầu; và kết hợp quy hoạch động.

\begin{enumerate}
  \item Chỉ thêm phương pháp cắt tỉa cải tiến, tức chương trình 2, không cải
  thiện nhiều hiệu suất thời gian. Lý do chính là độ chính xác của hàm ước
  lượng mới chỉ tăng nhẹ, trong khi hiệu suất thời gian của chính nó lại
  thấp, phần nào triệt tiêu lợi ích do ước lượng chính xác hơn đem lại.

  \item Phương pháp định cận ban đầu của bài này, tức chương trình 3, có độ
  chính xác rất cao, thường trực tiếp thu được lời giải tối ưu, nên có thể
  cải thiện hiệu suất khá rõ rệt. Trong bốn chương trình đầu, chỉ nó hoàn
  thành mọi dữ liệu trong bảng trong vòng một phút, đặc biệt hiệu quả khi
  \(n=1023\). Nhưng nếu độ chính xác của định cận ban đầu giảm một chút, hiệu
  quả tối ưu hóa sẽ trở nên rất không rõ, như trường hợp \(n=719\).

  \item Sau khi thêm quy hoạch động, chương trình 4 thực chất làm quá trình
  tìm kiếm chủ yếu tập trung dưới \(\lfloor n/2\rfloor\), tránh ở mức đáng kể
  việc tìm kiếm trong đoạn từ \(\lfloor n/2\rfloor\) đến \(n\), nơi hàm ước
  lượng gần như không có tác dụng. Vì vậy chương trình 4 gần như tốt hơn
  chương trình 1 khá rõ ở mọi điểm kiểm thử.
\end{enumerate}

Tiếp theo xét chương trình 5. Nhờ dùng định cận ban đầu, trước khi tìm kiếm
đã có một cận dưới độ ưu tương đối chính xác, nên hàm ước lượng mới có độ
chính xác cao hơn cũng có xác suất phát huy tác dụng cao hơn; do đó hiệu quả
tối ưu hóa khá rõ. Cũng như chương trình 3, khi định cận ban đầu không thật
chính xác, như \(n=719\), hiệu quả tối ưu hóa giảm mạnh.

Trong bảy chương trình, tốt nhất đương nhiên là chương trình 6. Vì tổng hợp
nhiều phương pháp tối ưu hóa và để chúng bổ sung ưu thế cho nhau, nó đạt hiệu
quả tối ưu hóa ổn định nhất.

Cuối cùng, ta còn thấy rằng trong trường hợp đặc thù của bài này, tìm kiếm
quay lui theo chiều sâu sau khi đã tối ưu hóa đầy đủ rõ ràng tốt hơn thuật
toán \(A^\ast\) cả về thời gian lẫn không gian. Chương trình 7 được biên dịch
ở chế độ bảo vệ và ngoài \(A^\ast\) còn kết hợp nhánh cận; nếu không, hiệu
suất thời gian và không gian còn kém hơn.

\subsection{Phụ lục chương trình}

\paragraph{Chương trình ``Betsy đi du lịch'' dùng hai kiểu cắt tỉa.}
\begingroup\small
\begin{verbatim}
{$A+,B-,D-,E+,F-,G+,I-,L-,N+,O-,P-,Q-,R-,S-,T-,V+,X+,Y-}
{$M 65520,0,655360}
program Betsy;{Vi du 1 cua bai viet doi tuyen tap huan IOI'99: Betsy di du lich}
{Ghi chu:
        1. De thuan tien do thoi gian, chuong trinh doc du lieu tu dong lenh.
        2. De xu ly tien hon, chuong trinh them mot vong o danh dau ngoai ban do.
        3. Cac doan tuong doi doc lap ve logic duoc tach bang dong trong.
}
const
    max=7;{Quy mo du lieu lon nhat ma chuong trinh nay xu ly duoc}
    delta:array[1..4,1..2]of shortint=((-1,0),(0,1),(1,0),(0,-1));{Do lech huong}


var
    map:array[0..max+1,0..max+1]of integer;
    {Ban do dung de danh dau lo trinh cua Betsy:
         vi tri chua di qua duoc danh dau -1,
         vong o danh dau ben ngoai duoc danh dau 0,
         cac o khac ghi so buoc khi Betsy den o do
    }
    left:array[0..max+1,0..max+1]of shortint;
    {Mang danh dau: ghi so o ke canh chua duoc Betsy di qua quanh moi o}
    n:1..max;{Canh ban do}
    n2:integer;{N*N}
    s:longint;{Bo dem, ghi tong so lo trinh di chuyen cua Betsy}


procedure init;{Doc du lieu va khoi tao}
    var
        i:integer;{Bien lap}
        temp:integer;{Bien tam}
    begin
        val(paramstr(1),n,temp);{Doc n tu dong lenh}
        n2:=n*n;


 fillchar(map,sizeof(map),255);
 for i:=0 to n+1 do begin
  map[0,i]:=0;
  map[i,0]:=0;
  map[n+1,i]:=0;
  map[i,n+1]:=0;
 end;
 map[1,1]:=1;
 {Doan tren khoi tao ban do}


 fillchar(left,sizeof(left),4);
 for i:=2 to n-1 do begin
  left[1,i]:=3;left[n,i]:=3;
  left[i,1]:=3;left[i,n]:=3;
 end;
 left[1,1]:=2;left[1,n]:=2;left[n,1]:=3;left[n,n]:=2;
 dec(left[1,2]);dec(left[2,1]);
 {Doan tren khoi tao mang danh dau}


 s:=0;{Xoa bo dem}
end;


procedure change(x,y,dt:integer);{Cong dt vao cac gia tri left cua 4 o ke canh (x,y)}
{Khi dat danh dau thi dt=-1; khi quay lui thi dt=1}
var
 k:integer;{Bien lap}
 a,b:integer;{Toa do tam}
begin
 for k:=1 to 4 do begin
  a:=x+delta[k,1];b:=y+delta[k,2];
  inc(left[a,b],dt);
 end;
end;


procedure expand(step,x,y:integer);{Thu tuc tim kiem chinh: tim buoc step, mo rong tu (x,y)}
var
  nx,ny:integer;{Vi tri moi cua Betsy}
  dir:byte;{Huong di chuyen cua Betsy}


 function able(x,y:integer):boolean;{Xet Betsy co the vao (x,y) hay khong}
  begin
    able:=not((map[x,y]<>-1)or((step<>n2)and(x=n)and(y=1)));
  end;


 function cut1:boolean;{Phan doan cat tia 1}
  var
    i:integer;{Bien lap}
    a,b:integer;{Toa do tam}
  begin
    for i:=1 to 4 do begin
     a:=x+delta[i,1];b:=y+delta[i,2];
     if (map[a,b]=-1)and((a<>nx)or(b<>ny))and(left[a,b]<=1) then begin
         cut1:=true;
         exit;
     end;
    end;
    cut1:=false;
  end;


 function cut2:boolean;{Phan doan cat tia 2}
  var
    d1,d2:integer;{Hai huong "trai" va "phai" so voi huong di hien tai}
    fx,fy:integer;{Vi tri neu Betsy di them mot buoc theo huong cu tu vi tri hien tai}
  begin
    if (dir=2)or(dir=4) then begin
     d1:=1;d2:=3;
    end
    else begin
     d1:=2;d2:=4;
    end;
    fx:=nx+delta[dir,1];fy:=ny+delta[dir,2];


  if (map[fx,fy]<>-1)and(map[nx+delta[d1,1],ny+delta[d1,2]]=-1)
  and(map[nx+delta[d2,1],ny+delta[d2,2]]=-1) then begin
   cut2:=true;
   exit;
  end;
  if (map[fx+delta[d1,1],fy+delta[d1,2]]<>-1)and(map[nx+delta[d1,1],ny+delta[d1,2]]=-1)
  and(map[fx,fy]=-1) then begin
   cut2:=true;
   exit;
  end;
  if (map[fx+delta[d2,1],fy+delta[d2,2]]<>-1)and(map[nx+delta[d2,1],ny+delta[d2,2]]=-1)
  and(map[fx,fy]=-1) then begin
   cut2:=true;
   exit;
  end;
  {Ba dieu kien tren tuong ung voi ba truong hop trong hinh nguyen ly cat tia}


  cut2:=false;
 end;


begin{Expand}
 if (step>n2)and(x=n)and(y=1) then begin{Tim den tang day, tang bo dem len 1}
  inc(s);
  exit;
 end;
 for dir:=1 to 4 do begin{Thu lan luot 4 huong di chuyen}
  nx:=x+delta[dir,1];ny:=y+delta[dir,2];
  if able(nx,ny) then begin
   if cut1 then continue;{Goi phan doan cat tia 1}
   if cut2 then continue;{Goi phan doan cat tia 2}
   change(nx,ny,-1);
   map[nx,ny]:=step;
   expand(step+1,nx,ny);{Goi de quy tang tim kiem tiep theo}
   map[nx,ny]:=-1;
   change(nx,ny,1);
  end;
     end;
    end;


begin{Chuong trinh chinh}
    init;{Khoi tao}
    expand(2,1,1);{Goi qua trinh tim kiem}
    writeln('The number of tours is ',s);{Xuat ket qua}
end.
\end{verbatim}
\endgroup

\paragraph{Chương trình ``số lần nhân ít nhất'', tức chương trình 6 trong phụ lục.}
\begingroup\small
\begin{verbatim}
{$A+,B-,D-,E+,F-,G+,I-,L-,N+,O-,P-,Q-,R-,S-,T-,V+,X+,Y-}
{$M 65520,0,655360}
program LeastMultiply;{Vi du 2 cua bai viet doi tuyen tap huan IOI'99: so lan nhan it nhat}
{Ghi chu:
     1. De thuan tien do thoi gian, chuong trinh doc n tu dong lenh.
     2. Sau khi ket thuc, chuong trinh ghi cach thuc hien phep nhan va tong so phep nhan vao output.txt.
     3. Trong qua trinh tim kiem, do co ket hop quy hoach dong, chuong trinh co the thu duoc do dai day
loi giai toi uu khi chua tim den tang day. Luc do chua co day so mu toi uu day du, nen sau khi tim kiem
ket thuc, chuong trinh goi formkeep de sinh day cau tao day du trong keep.
     4. Qua trinh tim kiem sinh day so mu toi uu, chua truc tiep dua ra cach thuc hien cac phep nhan,
nen thu tuc output chuyen doi no khi xuat ket qua.
     5. De nang cao hieu suat thoi gian, chuong trinh co vai toi uu hoa chi tiet; xem chu thich trong code.
     6. Cac doan tuong doi doc lap ve logic duoc tach bang dong trong.
}


const
    max=20;{Do dai lon nhat cua day}
    maxr=2000;{Do dai lon nhat cua mang dem quy hoach dong; n dau vao khong duoc vuot qua 2*maxr}
    mp=100;{Can tren cua cac so co the xu ly truc tiep bang tim kiem khi tien xu ly uoc luong}
    power2:array[0..12]of integer= (1,2,4,8,16,32,64,128,256,512,1024,2048,4096); {Luy thua cua 2}


type
    atype=array[0..max]of integer;{Kieu mang ghi day so mu duoc cau tao; phan tu 0 ghi do dai day}


var
    n:integer;{So dich doc vao}
    time:array[0..maxr]of integer;
    {Mang dem quy hoach dong: time[i] la so so hang toi thieu can dung de tao so i tren co so day hien tai}
    range:integer;{Pham vi xu ly bang quy hoach dong: range=[n/2]}
    a:atype;{Day ghi trong luc tim kiem}
    kp:atype;{Mang tam ghi day ket qua khi tien xu ly uoc can}
    keep:atype;{Mang ghi day ket qua toi uu}
    best:integer;{Loi giai toi uu hien tai}
    f:text;{Tep xuat}


procedure init;{Khoi tao}
    var
     temp:integer;{Bien tam}
    begin
     val(paramstr(1),n,temp);{Doc n tu dong lenh}
     keep[0]:=1;
     keep[1]:=1;
     best:=maxint;{Gia tri dau cua do dai day toi uu}
     assign(f,'output.txt');{Noi voi tep xuat}
    end;


procedure search(n:integer;var best:integer;var keep:atype);{Qua trinh tim kiem chinh}
{Truoc khi tim kiem, muc tieu tim kiem la n;
            best luu can duoi do uu da co truoc khi tim kiem;
            keep luu day so mu toi uu ung voi can duoi do uu ban dau.
Sau khi tim kiem ket thuc, best la do dai day so mu toi uu duoc cau tao,
tuc so lan nhan it nhat cong 1; keep la day so mu duoc cau tao.
}
    var
     kn:integer;{So lan xuat hien thua so 2 trong n}
     i:integer;{Bien lap}


 function getk(num:integer):integer;{Tinh so lan xuat hien thua so 2 trong num, tuc k(num) trong bai}
  var
    i:integer;{Bien lap}
  begin
    i:=0;
    while not odd(num) do begin
     num:=num shr 1;
     inc(i);
    end;
    getk:=i-1;
  end;


 procedure find(step:integer);{Qua trinh tim kiem de quy}
  var
    i:integer;{Bien lap}
    k:integer;{Can tren pham vi lap cua tang nay}


  function ok(num:integer):boolean;{Xet so num co the duoc cau tao tai thoi diem hien tai hay khong}
  {De nang cao hieu suat, o day dung ket qua quy hoach dong}
    var
     i,j:integer;{Bien lap}
    begin
     if num<=range then begin{So num can xet nam trong [n/2]}
         ok:=(time[num]=2);{Dung truc tiep dieu kien so so hang toi thieu co bang 2 hay khong}
         exit;
     end;
     for i:=step-1 downto 1 do begin
         if a[i]+a[i]<num then break;
         if time[num-a[i]]=1 then begin
         {time[t]=1 cho biet t da xuat hien trong day hien co, tranh phai dung vong lap de xet}
          ok:=true;
          exit;
         end;
     end;
     ok:=false;
    end;


 procedure evltime;{Thu tuc con quy hoach dong}
  var
   i,j:integer;{Bien lap}
   p:integer;{Can tren cua lan truy hoi quy hoach dong nay}
  begin
   p:=k;
   if p>range then p:=range;
   for i:=a[step-1]+1 to p do begin
    time[i]:=time[i-a[step-1]]+1;{Gan gia tri dau cho ham muc tieu}
    for j:=step-2 downto 2 do{Liet ke quyet dinh}
        if time[i-a[j]]+1<time[i] then time[i]:=time[i-a[j]]+1;
   end;
  end;


 function h(num:integer):byte;{Ham uoc luong don gian ban dau h}
  var
   i:integer;{Bien dem vong lap}
  begin
   i:=0;
   while num<n do begin
    num:=num shl 1;
    inc(i);
   end;
   h:=i;
  end;


 function cut1:boolean;{Phan doan cat tia ban dau, goi truc tiep ham uoc luong h}
  begin
   if h(a[step])+step>=best then cut1:=true else cut1:=false;
  end;


 function cut2:boolean;{Phan doan cat tia dung ham uoc luong cai tien}
  var
   pt:integer;{Gia tri k(a[step]), tuc so mu cua 2 trong a[step]}
   rest:integer;{So trong day moi thoa k(rest)=k(n)}
     i:integer;{Bien dem vong lap}
    begin
     if h(a[step]+a[step-1])+step+1<best then begin
     {Neu buoc dau cua day moi chon a[step]+a[step-1] thay vi 2*a[step] thi chac chan se bi cat tia;
      day la dieu kien can de ham uoc luong moi phat huy tac dung.
     }
         cut2:=false;
         exit;
     end;


     pt:=getk(a[step]);{Tinh k(a[step])}
     if pt>kn then rest:=a[step-1]
     else
         if pt=kn then rest:=a[step]
         else rest:=maxint;
           {Gan gia tri dau cho rest de phong khi moi so trong day moi co k nho hon k(n)}


     i:=0;
     repeat
         a[step+i+1]:=a[step+i] shl 1;
         inc(i);
         if pt+i=kn then rest:=a[step+i];
     until a[step+i]>n;
     dec(i);
     {Doan tren la qua trinh cau tao day}


     if (n-a[step+i]>rest)and(step+i+2>=best) then cut2:=true else cut2:=false;{Phan doan cat tia}
    end;


  begin{Find}
    if a[step-1]+a[step-1]>=n then begin
    {So lon nhat hien tai trong day da vuot qua [n/2], nen dung truc tiep ket qua quy hoach dong}
     if time[n-a[step-1]]+step-1<best then begin{Xac dinh co loi giai}
         best:=time[n-a[step-1]]+step-1;
         keep:=a;
     end;
      exit;
     end;


     k:=a[step-1]+a[step-1];{Tinh can tren pham vi lua chon cua a[step]}
     evltime;{Goi thu tuc con quy hoach dong}


     inc(a[0]);
     for i:=k downto a[step-1]+1 do
      if ok(i) then begin
          if i<=range then time[i]:=1;
          a[step]:=i;
          if cut1 then break;{Goi phan doan cat tia 1}
          if cut2 then continue;{Goi phan doan cat tia 2}
          find(step+1);{Goi de quy tang tim kiem tiep theo}
      end;
     dec(a[0]);
    end;


procedure formkeep;{Sinh day toi uu keep}
{Trong luc tim kiem, neu a[step]>[n/2] thi dung truc tiep ket qua quy hoach dong,
nen phan cuoi cua day ket qua toi uu thuc ra chua duoc tao ra. Qua trinh nay
bo sung phan do. O day van can de quy va quay lui, thuc chat la mot tim kiem
quy mo nho; tuy nhien, vi ket qua quy hoach dong bieu thi viec chon so it nhat
cac so tu day hien co de cong thanh n, nen cac so nay co the duoc sap thu tu.
}
    var
     found:boolean;{Co danh dau da tim thay day toi uu}


    procedure check(from,left:integer);{Quay lui sinh phan cuoi cua day toi uu}
      {from la can tren cua vong lap trong tang hien tai, dung de sap thu tu;
       left la phan con lai can thu duoc bang phep cong}
     var
      i:integer;{Bien lap}
     begin
      if keep[0]=best then begin
          if left=0 then found:=true;{Tim thay day toi uu}
         exit;
     end;


     inc(keep[0]);
     for i:=from downto 1 do{Liet ke cac so trong day}
         if keep[i]<=left then begin
          keep[keep[0]]:=keep[keep[0]-1]+keep[i];
          check(i,left-keep[i]);{Goi tang tim kiem tiep, nhung so duoc dung khong vuot qua keep[i]}
          if found then exit;
         end;
     dec(keep[0]);
    end;


  begin{FromKeep}
    found:=false;
    check(keep[0],n-keep[keep[0]]);{Goi qua trinh tim kiem}
  end;


 begin{Search}
  kn:=getk(n);{Tinh k(n)}
  time[0]:=0;
  time[1]:=1;
  a[0]:=1;
  a[1]:=1;
  range:=n div 2;
  {Doan tren khoi tao truoc khi tim kiem}


  find(2);{Goi qua trinh tim kiem}
  formkeep;{Sau khi tim kiem ket thuc, sinh day cau tao day du trong keep}
 end;


function guess(n:integer):integer;{Tien xu ly uoc can bang de quy}
{Ghi chu:
   1. Sau khi chuong trinh con guess ket thuc, gia tri tra ve la ket qua uoc luong cho n;
      dong thoi, mang keep nhan duoc day tuong ung.
   2. De uoc luong cang chinh xac cang tot, guess xet dong thoi hai chien luoc phan tich,
      tuc la dung cau truc quay lui.
   3. Vi moi lan goi de quy guess deu phai ghi ket qua cuoi cung vao keep, moi lan guess
      chi thao tac mot phan cua mang keep va can den cac mang tam kp, kpt.
}
    var
     num:integer;{Phan con lai sau khi tach het cac thua so 2 trong n}
     pfact:integer;{Bien bieu thi thua so nguyen to cua num}
     best:integer;{Ghi gia tri uoc luong tra ve khi goi xuong duoi}
     b2:integer;{Ghi so thua so 2 trong num}
     s:integer;{So buoc uoc luong cho n}
     i:integer;{Bien lap}
     sq:integer;{Phan nguyen can bac hai cua num, lam can tren khi phan tich thua so}
     s1,k2,k:integer;{Bien tam}
     kpt:atype;{Tam ghi day ket qua toi uu}


    begin
     num:=n;
     b2:=0;
     while not odd(num) do begin
      num:=num div 2;
      inc(b2);
      inc(keep[0]);
      keep[keep[0]]:=power2[b2];
     end;
     s:=b2+1;
     k2:=keep[0];
     {Doan tren tach het cac thua so 2 trong n, phan con lai la num}


     if num<=mp then begin{Neu num du nho, khong vuot qua mp, thi goi truc tiep qua trinh tim kiem}
      best:=maxint;
      search(num,best,kp);{Goi tim kiem cho n, day thu duoc luu trong kp}
      inc(s,best-1);
      for i:=2 to best do keep[i+keep[0]-1]:=kp[i];
      inc(keep[0],best-1);
     end
     else begin{Xu ly bang phuong phap uoc luong}
      k:=keep[0];


    best:=guess(num-1);{Goi de quy guess}
    keep[keep[0]+1]:=keep[keep[0]]+1;
    inc(keep[0]);
    inc(s,best);
    {Doan tren la chien luoc uoc luong thu nhat: tru 1 khoi num roi uoc luong num-1}


    sq:=trunc(sqrt(num));
    pfact:=3;
    while (pfact<=sq)and(num mod pfact>0) do inc(pfact,2);
    if pfact<=sq then begin
        kpt:=keep;
        keep[0]:=k2;
        s1:=b2+1;
        {Doan tren khoi tao truoc khi dung chien luoc thu hai}


        k:=keep[0];
        best:=guess(pfact);
        inc(s1,best-1);
        {Doan tren goi de quy guess de uoc luong thua so nguyen to pfact}


        k:=keep[0];
        best:=guess(num div pfact);
        for i:=k+1 to keep[0] do keep[i]:=pfact*keep[i];
        inc(s1,best-1);
        {Doan tren uoc luong [num/pfact]}


        if s1<s then s:=s1 else keep:=kpt;{So sanh hai chien luoc}
    end;
    {Doan tren la chien luoc uoc luong thu hai: tach mot thua so nguyen to nho cua num roi xu ly.
     Ket qua uoc luong duoc luu trong s1; kpt tam luu day do chien luoc thu nhat tao ra.
    }
  end;


  for i:=k2+1 to keep[0] do keep[i]:=keep[i]*power2[b2];
  guess:=s;{Tra ve ket qua uoc luong}
 end;


procedure output;{Xuat ket qua}
 var
  i,j,k:integer;{Bien lap}
 begin
  rewrite(f);
  writeln(f,'X1');
  for i:=2 to best do begin
   for j:=1 to i-1 do begin
       for k:=j to i-1 do
        if keep[j]+keep[k]=keep[i] then break;
       if keep[j]+keep[k]=keep[i] then break;
   end;
   writeln(f,'X',i,'=X',j,'*X',k);
  end;
  writeln(f,'Times of multiplies=',best-1);
  close(f);
 end;


begin{Chuong trinh chinh}
 init;{Doc du lieu va khoi tao}
 best:=guess(n);{Goi tien xu ly uoc can}
 search(n,best,keep);{Goi qua trinh tim kiem chinh}
 output;{Xuat ket qua}
end.
\end{verbatim}
\endgroup

\end{document}
